% OC Core 1.3 — Cross-level structures master module
% This file orchestrates all cross-K linking sections.

\section{Межуровневые структуры}
\label{sec:crossk_master}

\subsection{Назначение слоя Cross-K}
Слой Cross-K формализует все структуры, переходы и ограничения,
которые \emph{связывают различные уровни континуумов $K_0 \rightarrow K_{12}$}.
Пока каждый уровень $K$ имеет собственную онтологию (оси, потенциалы,
пороги, потоки, циклы, области $\Omega(K)$, границы $\partial\Omega(K)$ и
меру континуума $k(K)$), в Core множество явлений требуют
\textbf{отношений между уровнями}. Эти отношения не являются ни
необязательными, ни внешними: они составляют неотъемлемую архитектуру
динамического континуума.

Поэтому слой Cross-K выполняет:
\begin{itemize}
  \item определение \textbf{структур отображения} между уровнями;
  \item задание \textbf{операторов перехода} и их ограничений;
  \item введение \textbf{условий согласованности}, гарантирующих математическую
        и концептуальную корректность эволюции между смежными уровнями;
  \item определение \textbf{правил глобального ландшафта}, описывающих, как
        весь стек уровней $K$ формирует согласованный мета-континуум.
\end{itemize}

\subsection{Определение структур Cross-K}
 \emph{Структурой Cross-K} является любой формальный объект, для которого:
\begin{enumerate}
  \item затрагивает по крайней мере два уровня $K_i, K_j$ с $i \neq j$;
  \item устанавливает отношение $R_{ij}$ между их компонентами
        $(A,P,\Theta,J,C,\Omega,\partial\Omega,k)$;
  \item удовлетворяет общему правилу совместимости:
        \[
        R_{ij} : K_i \leftrightarrow K_j
        \]
        при сохранении непрерывности и разрешении корректных переходов в рамках
        глобального оператора $E$.
\end{enumerate}

Следовательно, в Cross-K входят:
\begin{itemize}
  \item \textbf{карты переходов} (размерностные, структурные, энергетические);
  \item \textbf{отношения наследования} (оси, потенциалы, пороги);
  \item \textbf{преобразования границ} между $\partial\Omega(K_i)$ и
        $\partial\Omega(K_j)$;
  \item \textbf{проекции и подъёмы циклов} (отображение циклов с одного уровня
        на другой);
  \item \textbf{глобальные ограничения}, предотвращающие нарушение соседними
        уровнями правил друг друга.
\end{itemize}

\subsection{Роль данного мастер-файла}
Этот мастер-модуль намеренно остаётся \textbf{тонким слоем оркестровки}.
Он не хранит тяжёлый теоретический материал. Вместо этого он:
\begin{itemize}
  \item вводит понятие структур Cross-K,
  \item даёт глобальное определение и назначение,
  \item и собирает все детальные межуровневые анализы через \verb|\input|.
\end{itemize}

Индивидуальные секции переходов организованы как:
\begin{itemize}
  \item глобальный ландшафт (всесторонняя перспектива),
  \item локальные переходы ($K_i \rightarrow K_{i+1}$),
  \item специальные двунаправленные или несмежные отображения, где это требуется.
\end{itemize}

\subsection{Структура модуля}
\noindent По этому модулю оркестрируются следующие файлы:
\begin{itemize}
  \item \verb|crossk_global_landscape.tex| — глобальная многоуровневая геометрия;
  \item \verb|crossk_k0_k1.tex|, \verb|crossk_k1_k2.tex|, \dots —
        блоки переходов между соседними уровнями;
  \item \verb|crossk_k10_k11.tex|, \verb|crossk_k11_k12.tex| — верхние
        переходы и мета-структуры.
\end{itemize}

\subsection{Включение детальных секций}
\noindent Детальные анализы Cross-K можно включать как:
\begin{verbatim}
\input{content/crossk/crossk_global_landscape.tex}
\input{content/crossk/crossk_k0_k1.tex}
\input{content/crossk/crossk_k1_k2.tex}
...
\input{content/crossk/crossk_k11_k12.tex}
\end{verbatim}

\noindent Таким образом, данный файл определяет \textbf{концептуальный контейнер}
для домена Cross-K OC Core 1.3.
